\scanpage{18}

un vectoriel \inserted{$V$} de rg $2$ sur $\mathbb F_2$, la donnée d'un $a\in\deleted{\unsure{\mathbb F_2|\mathbb F_2^*}}V^*$ revient à la donnée d'une partie de $V$ qui soit une base (savoir $V^*\setminus\{a\}$) ; donc la donnée d'un \inserted{couple $(V,\underline a)$ d'un} vectoriel $V$ de rg $2$ sur $\mathbb F_2$, et d'un élément \inserted{$a$} de $V^*$, revient à la donnée d'un ens. $C$ à deux éléments, par $V\simeq\mathbb F_2^C$, $a$ l'élément diagonal de $\mathbb F_2^C$. Ainsi, la donnée d'un carré \deleted{\unsure{rev}}$Q$ \deleted{\unsure{via les ens.}} (via l'ens des sommets) équivaut à celle d'un couple $(C,S)$, où $C$ \deleted{\unsure{$(\mathbb F_2)$}} est un ens à deux éléments, et $S$ est un torseur sous $\mathbb F_2^C$. Mais la donnée d'un tel torseur revient à celle d'une famille de \inserted{$\mathbb F_2$-}torseurs indexée par $C$, i.e. d'une application de degré $2$
\[
A\longrightarrow C,
\]
% slanted
\margin{ici les \unsure{deux|dits|\deleted{deux}} ens $A_c$ $(c\in C)$ \unsure{apparaissent|se présentent} \deleted{\unsure{comme}} comme \uline{quotients} de $S$, par les relations d'équivalence définies resp. par \deleted{les} \deleted{deux} \inserted{l'élément} \deleted{\unsure{$c$}} \inserted{$c'=c+a$} de $V^*$ (considérés comme involutions ss pt fixes de $S$) Mais je préfère ici prendre la convention inverse, en \deleted{identifiant} intervertissant $A_c$ et $A_{c'}$.}

$A$ a quatre éléments, et est muni d'une involution sans pts fixes (qui donne le passage au quotient vers $C$), par la formule
\[
S\simeq\prod_{c\in C}\deleted{\unsure{\mathbb F_2}} A_c=\text{Ens des sections de } A\text{ sur } C
\]
qui donne en même temps un plongement
\[
S\hookrightarrow\mathscr P_2(A)\unsure{.}
\]
d'où la relation d'incidence entre $S$ et $A$\ldots

% archive sheet number: 17

